\addcontentsline{toc}{subsubsection}{Solution 47. The Three-Variable Quadratic Trap}

\begin{dlsummary}[title={\bfseries Solution 47. The Three-Variable Quadratic Trap},breakable]

\begin{proof}
The unique solution is
\[
\boxed{f(x)=x^2\qquad(x\in\mathbb R).}
\]

\paragraph{1. Force evenness and exact doubling.}

Condition (2) at \(x=0\), together with nonnegativity, gives \(f(0)=0\).

In (1), substitute \((x,y,z)=(t,t,-t)\):
\[
f(2t)\le3f(t)+f(-t).
\]
Combining this with (2) gives \(f(t)\le f(-t)\). Replacing \(t\) by \(-t\) proves
\[
f(-t)=f(t).
\]
The same inequalities now yield
\[
\boxed{f(2t)=4f(t).} \tag{3}
\]

\paragraph{2. Make the inequality reverse itself.}

Substitute \(z=-y\) in (1). Evenness gives
\[
f(x+y)+f(x-y)\le2f(x)+2f(y). \tag{4}
\]

Now apply (4) to \(x+y\) and \(x-y\):
\[
f(2x)+f(2y)\le2f(x+y)+2f(x-y).
\]
Using (3), this is the reverse of (4). Consequently,
\[
\boxed{f(x+y)+f(x-y)=2f(x)+2f(y).} \tag{5}
\]

\paragraph{3. Determine a dense set of values.}

Taking \(y=1\) in (5),
\[
f(n+1)+f(n-1)=2f(n)+2.
\]
Starting from \(f(0)=0\) and \(f(1)=1\), induction gives
\[
f(n)=n^2\qquad(n\in\mathbb Z).
\]
Exact doubling then gives
\[
f\!\left(\frac{m}{2^k}\right)=\left(\frac{m}{2^k}\right)^2
\]
for every integer \(m\) and nonnegative integer \(k\).

\paragraph{4. Derive continuity instead of assuming it.}

Equations (3) and (5) imply midpoint convexity:
\[
f\!\left(\frac{u+v}{2}\right)
\le\frac{f(u)+f(v)}2.
\]
Iterating this gives the corresponding inequality for dyadic weights.

Fix \(x\), and let \(N\) be a power of two. If \(|h|\le1/N\), then evenness and the assumed unit-interval bound give \(f(Nh)\le1\). By (5) and nonnegativity,
\[
f(x+Nh)\le2f(x)+2f(Nh)\le2f(x)+2.
\]
Using the dyadic weight \(1/N\),
\[
f(x+h)\le
\left(1-\frac1N\right)f(x)+\frac1N f(x+Nh)
\le f(x)+\frac{f(x)+2}{N}.
\]
The same upper bound holds for \(f(x-h)\). Since midpoint convexity gives
\[
2f(x)\le f(x+h)+f(x-h),
\]
we obtain
\[
|f(x+h)-f(x)|\le\frac{f(x)+2}{N}.
\]
Thus \(f\) is continuous everywhere.

Approximating any real number by dyadic rationals now proves \(f(x)=x^2\).

\paragraph{5. Verify.}

For \(f(x)=x^2\), expansion shows that (1) holds with equality. All other conditions are immediate.
\end{proof}

\paragraph{Takeaways.}
\begin{itemize}
\item **Exploit sign changes:** The substitution \((t,t,-t)\) creates evenness from a one-sided inequality.
\item **Change coordinates:** Replacing \((x,y)\) by \((x+y,x-y)\) reverses the parallelogram inequality.
\item **Local bounds can replace continuity:** A bound on one interval, combined with the functional structure, controls the function everywhere.
\item **Do not stop at rational values:** Extending a classification to all real inputs requires a justified regularity argument.
\end{itemize}

\end{dlsummary}
